\cleardoublepage

\begin{center}
    {\LARGE\bfseries Sommario}
\end{center}

\vspace{1cm}

Il calcolo delle proprietà macroscopiche di un sistema fisico a partire dalla sua descrizione microscopica richiede, nell'ambito della meccanica statistica, la valutazione di valori di aspettazione rispetto alla distribuzione di Boltzmann. Poiché lo spazio delle configurazioni ha dimensione pari al numero di gradi di libertà del sistema, tale calcolo è generalmente inaccessibile per via analitica o mediante integrazione numerica, ed è affrontato tramite metodi Monte Carlo, che sostituiscono il calcolo esatto con una stima statistica basata sul campionamento.
 
Un metodo di campionamento a cui si ricorre è \emph{Markov Chain Monte Carlo} (MCMC), tipicamente utilizzato con proposte locali. Per quanto generale, questo approccio mostra un limite intrinseco quando la distribuzione target è multimodale: ovvero quando presenta modi di alta densità separati da regioni a bassa densità di probabilità. In questo caso la catena di Markov può rimanere intrappolata in un singolo modo per tempi computazionali proibitivi, compromettendo il corretto capionamento (fenomeno noto come \emph{intrappolamento nei modi}).
 
Questa tesi esplora una strategia alternativa basata sui \emph{Normalizing Flows} (NF), una classe di modelli generativi costruiti tramite trasformazioni invertibili e differenziabili. Un NF addestrato ad approssimare la distribuzione di Boltzmann può essere impiegato in due modi: come motore di campionamento diretto, le cui stime vengono poi corrette tramite \emph{importance sampling}; oppure come distribuzione di proposta globale all'interno di in un metodo ibrido (Flow-MCMC), che unisce la capacità di esplorazione globale del Normalizing Flow alla correttezza asintotica garantita da Markov Chain Monte Carlo.

I tre metodi - MCMC, NF e Flow-MCMC - vengono confrontati empiricamente su un sistema di test a potenziale doppio pozzo, scelto per la bimodalità che induce alla distribuzione di Boltzmann. Si quantifica l'efficienza del campionamento al variare della temperatura e del numero di campioni. I risultati mostrano che, a basse temperature, l'MCMC locale rimane sistematicamente confinato in un solo pozzo del potenziale, producendo stime distorte. NF e Flow-MCMC, al contrario, mantengono un'esplorazione bilanciata dei due modi e una convergenza più rapida su tutto l'intervallo di temperature considerato, confermando l'assenza di intrappolamento nei modi e dunque il vantaggio pratico delle mosse globali basate sui Normalizing Flows rispetto al campionamento markoviano locale in presenza di distribuzioni di probabilità multimodali.